\section{Revisão de literatura}

Este capítulo é construído em duas camadas. A primeira, apresentada na
Seção~\ref{sec:leituras}, é uma proposta de leitura: cinco artigos que
sustentam, em conjunto, as principais escolhas relativas à medida de exposição
e à estratégia de identificação adotadas no Capítulo~\ref{sec:metodo}; decisões
como pareamento, ponderação e inferência têm justificativa própria naquele
capítulo. Não são cinco textos escolhidos por proximidade
temática genérica com o tema ``inteligência artificial e trabalho'', e sim os
cinco textos de que este trabalho depende para justificar, respectivamente, a
medida de tratamento, a medida alternativa de exposição, o que se sabe sobre o
efeito de produtividade da ferramenta, o que se sabe sobre o efeito agregado no
mercado de trabalho e o desenho de identificação por exposição diferencial. A
segunda camada, nas seções seguintes, é a revisão propriamente dita, construída a partir dessas leituras e ampliada com a literatura brasileira de
mobilidade ocupacional e informalidade.

\subsection{Proposta de leitura}
\label{sec:leituras}

O Quadro~\ref{qua:leituras} lista as cinco leituras propostas, na ordem em que
convém lê-las. A ordem não é cronológica: começa pela medida que este trabalho
usa como tratamento, passa pela medida concorrente, desce ao nível da firma e
do trabalhador e termina no desenho econométrico que serve de modelo ao
Capítulo~\ref{sec:metodo}.

\begin{quadro}[htb]
  \caption{Proposta de leitura, na ordem sugerida}
  \label{qua:leituras}
  \centering
  \espacosimples\idpcorpodez
  \begin{tabular}{@{}p{0.5cm}p{4.3cm}p{9.4cm}@{}}
    \toprule
    \textbf{\#} & \textbf{Referência} & \textbf{Por que ler} \\
    \midrule
    1 & \citeonline{felten2021}
      & Define o AIOE, que é literalmente a variável de tratamento deste
        trabalho. Sem esse texto não se sabe o que o coeficiente mede. \\
    \addlinespace
    2 & \citeonline{eloundou2024}
      & Medida concorrente de exposição, específica a modelos de linguagem.
        É a primeira robustez prevista no Capítulo~\ref{sec:melhorias}. \\
    \addlinespace
    3 & \citeonline{brynjolfsson2025}
      & Efeito estimado com adoção observada e implantação escalonada. Serve de
        contraste: mostra o que este trabalho não consegue observar. \\
    \addlinespace
    4 & \citeonline{humlum2025}
      & Desenho com exposição ocupacional e horizonte curto, com efeitos nulos
        e precisos. Referência para calibrar a expectativa de magnitude. \\
    \addlinespace
    5 & \citeonline{acemoglu2020}
      & Modelo canônico de identificação por exposição diferencial e leitura
        de deslocamento \emph{versus} recomposição. \\
    \bottomrule
  \end{tabular}
  \fonte{Elaboração própria (2026).}
\end{quadro}

As subseções a seguir retomam essas leituras na ordem do
Quadro~\ref{qua:leituras}, e todas seguem o mesmo percurso: sintetizam,
inicialmente, a contribuição do artigo e, na sequência, explicitam por que a
leitura importa para esta dissertação, seja pela decisão de desenho que
sustenta, seja pelo limite que ajuda a reconhecer. A coluna ``Por que ler'' do
quadro antecipa, em forma resumida, o argumento que cada subseção desenvolve.
Lidas em conjunto, as cinco leituras formam a base sobre a qual se apoiam as
seções seguintes deste capítulo e as escolhas metodológicas apresentadas no
Capítulo~\ref{sec:metodo}.

\subsubsection{Felten, Raj e Seamans (2021)}

\citeonline{felten2021} constroem o \textit{Artificial Intelligence
Occupational Exposure} (AIOE), índice que liga aplicações correntes de
inteligência artificial às habilidades exigidas por cada ocupação e, a partir
daí, produz um escore de exposição por ocupação da classificação ocupacional
americana. O índice é padronizado entre ocupações, de modo que uma unidade
corresponde aproximadamente a um desvio-padrão de exposição.

Trata-se do artigo mais importante da lista por uma razão simples: o AIOE é o
tratamento deste trabalho, e sua leitura responde às dúvidas que o leitor tende
a formular diante do índice. O AIOE mede a exposição potencial de cada ocupação
às aplicações de inteligência artificial, e não a adoção efetivamente observada
da tecnologia, seus efeitos sobre a produtividade ou o risco de substituição do
trabalho. A padronização entre ocupações, por sua vez, define a escala em que se
leem os coeficientes estimados no Capítulo~\ref{sec:resultados}, expressos por unidade do índice. Por fim, exposição elevada não deve ser tomada como
sinônimo de ocupação pior, pois o índice descreve a proximidade entre as
habilidades requeridas pela ocupação e as capacidades da tecnologia, e não a
qualidade do emprego. Essa distinção precisa estar clara antes de qualquer
interpretação distributiva dos resultados.

\subsubsection{Eloundou et al. (2024)}

\citeonline{eloundou2024} propõem uma medida de exposição construída
especificamente para modelos de linguagem de grande porte, avaliando tarefa a
tarefa se o acesso a um modelo reduziria substancialmente o tempo de execução.
Estimam que cerca de 80\% da força de trabalho americana teria ao menos 10\%
das tarefas afetadas e que aproximadamente 19\% teria ao menos metade delas.

A leitura é necessária por uma razão de robustez. O AIOE foi construído antes
da difusão dos modelos generativos e capta exposição à inteligência artificial
em sentido amplo; a medida de \citeonline{eloundou2024} é contemporânea ao
choque que este trabalho usa como marco temporal. Se as duas medidas produzem
gradientes de sinal e magnitude semelhantes, o resultado é mais crível; se
divergem, a divergência informa qual dimensão de exposição está de fato
correlacionada com as transições observadas. A substituição da medida de
exposição está registrada como primeira melhoria posterior na
Seção~\ref{sec:exposicoes}.

\subsubsection{Brynjolfsson, Li e Raymond (2025)}

\citeonline{brynjolfsson2025} acompanham a implantação escalonada de um
assistente conversacional em uma central de atendimento e medem ganho de
produtividade, com efeitos concentrados entre trabalhadores menos experientes.
A estratégia apoia-se na observação da adoção, na variação do momento de
implantação entre equipes e na medição do desfecho no nível do trabalhador;
como em todo desenho de implantação escalonada, a leitura causal depende de que
esse momento não esteja associado a tendências próprias das equipes.

O artigo entra na lista como contraste metodológico deliberado. Ele mostra o
que se consegue afirmar quando a adoção é observada e o que, por consequência,
não se consegue afirmar quando ela não é. Neste trabalho não há adoção
observada: há um índice ocupacional e um marco temporal nacional. A comparação com \citeonline{brynjolfsson2025} delimita, assim, o que os
coeficientes do Capítulo~\ref{sec:resultados} podem e não podem sustentar.

\subsubsection{Humlum e Vestergaard (2025)}

\citeonline{humlum2025} combinam levantamentos representativos de adoção de
assistentes conversacionais com registros administrativos vinculados de
empregador e empregado na Dinamarca. Estimam efeitos nulos e precisos sobre
rendimentos e horas, descartando efeitos superiores a poucos pontos percentuais
dois anos após a adoção, e encontram recomposição de tarefas e alguma troca de
ocupação sem variação líquida de horas ou remuneração. O texto circulou
inicialmente com o título \textit{Large Language Models, Small Labor Market
Effects}.

É o artigo cuja leitura mais muda a expectativa sobre a magnitude dos
resultados. O desenho compartilha elementos com o deste trabalho (diferenças em
diferenças, exposição ocupacional e horizonte curto), mas usa dados
administrativos muito mais ricos e examina desfechos e populações distintos. O
achado central, de que a margem que se move é a de composição de tarefas e de
troca de ocupação, e não a de emprego ou salário, é parcialmente compatível com
o padrão dos resultados do Capítulo~\ref{sec:resultados}, em que o coeficiente
estatisticamente mais forte é o de deslocamento no gradiente de exposição, e não
o de saída do emprego. Trata-se de compatibilidade parcial, e não de
confirmação, dadas as diferenças de desenho, população e desfecho.

\subsubsection{Acemoglu e Restrepo (2020)}

\citeonline{acemoglu2020} estimam o efeito da adoção de robôs industriais sobre
emprego e salários em mercados de trabalho locais americanos, usando a exposição
diferencial desses mercados à robotização. É o desenho canônico de identificação
por exposição diferencial e a referência padrão para a distinção entre efeito de
deslocamento e efeito de recomposição.

A leitura serve a dois propósitos. O primeiro é metodológico: o trabalho mostra
como se argumenta a favor de uma medida de exposição como fonte de variação
identificadora e, mais importante, quais objeções essa estratégia atrai:
tendências pré-existentes, choques setoriais concomitantes, seleção na medida de
exposição. Todas reaparecem no Capítulo~\ref{sec:metodo}. O segundo é conceitual:
a decomposição entre deslocamento e recomposição dá vocabulário para interpretar
por que a permanência ocupacional cai em todas as faixas de exposição no período
analisado, e não apenas nas mais expostas.

\subsection{Inteligência artificial, automação e mercado de trabalho}

A literatura econômica sobre automação e trabalho organiza-se em torno de uma
decomposição que \citeonline{acemoglu2019} formalizaram no quadro de tarefas.
Toda tecnologia que passa a executar tarefas antes desempenhadas por pessoas
produz, ao mesmo tempo, um \emph{efeito de deslocamento}, que reduz a demanda
por trabalho nas tarefas automatizadas, e um \emph{efeito de recomposição}, que
a amplia nas tarefas remanescentes, nas novas tarefas criadas pela própria
tecnologia e nos setores beneficiados pelo ganho de produtividade. O sinal do
efeito líquido não é, por isso, uma predição teórica: é uma questão empírica,
cuja resposta depende de qual dos dois efeitos domina em cada horizonte e em
cada mercado. A onda anterior de automação (robótica industrial e
informatização de rotinas) alcançou sobretudo tarefas codificáveis e
repetitivas, concentradas em ocupações de escolaridade intermediária, e a
evidência de \citeonline{acemoglu2020} para mercados de trabalho locais
americanos, construída sobre a exposição diferencial desses mercados à
robotização, indica predomínio do deslocamento nesse caso. Dessa literatura este
trabalho herda não apenas o vocabulário, mas a estratégia de identificação: usar
a exposição diferencial à tecnologia como fonte de variação, com as objeções que
ela atrai e que reaparecem no Capítulo~\ref{sec:metodo}.

A inteligência artificial generativa desloca essa fronteira. Diferentemente da
automação de rotinas, ela alcança tarefas cognitivas não rotineiras (redigir,
resumir, programar, atender), historicamente protegidas por exigirem
linguagem e julgamento, e concentradas em ocupações de escolaridade alta. A
predição de viés de rotina, que orientou a leitura distributiva da onda
anterior, deixa de valer automaticamente: não há razão \textit{a priori} para
que a exposição à inteligência artificial se distribua na mesma direção da
exposição à robotização, e é justamente por isso que a direção do gradiente
precisa ser estimada em vez de suposta. A evidência disponível até aqui, porém,
é mais contida do que o debate público sugere, e convém não tomá-la como bloco
homogêneo, porque os dois trabalhos de referência medem objetos distintos.
\citeonline{brynjolfsson2025} documentam ganhos de produtividade concentrados
entre trabalhadores menos experientes, em desenho com adoção observada no
interior de uma única firma: é evidência sobre desempenho individual na execução
da tarefa, não sobre emprego ou salário agregados. \citeonline{humlum2025}, com
registros administrativos vinculados e adoção declarada em escala nacional,
estimam efeitos nulos e precisos sobre rendimentos e horas dois anos após a
adoção, e observam recomposição de tarefas e alguma troca de ocupação: é
evidência sobre resultados agregados de mercado, não sobre produtividade
individual. As duas são complementares, e não redundantes: a primeira mostra
que a ferramenta entrega ganho onde é adotada; a segunda, que esse ganho ainda
não se converteu em deslocamento mensurável de renda ou de horas no horizonte
observado. Tomadas em conjunto, tornam plausível que a margem de ajuste mais
provável no curto prazo seja a de composição e realocação, e é essa expectativa,
e não a de destruição líquida de postos, que orienta a leitura dos
resultados no Capítulo~\ref{sec:resultados}.

\subsection{Medidas de exposição ocupacional à inteligência artificial}

As medidas de exposição ocupacional à inteligência artificial usadas na
literatura empírica dividem-se em duas famílias, e a diferença entre elas não é
de calibragem, mas de objeto. A primeira parte das \emph{habilidades}
requeridas pelas ocupações: o AIOE de \citeonline{felten2021} associa aplicações
de inteligência artificial às habilidades requeridas pelas ocupações e agrega
essa correspondência em um escore padronizado entre ocupações. Mede, portanto,
quanto do repertório de habilidades de uma ocupação está ao alcance da
tecnologia. A segunda parte das \emph{tarefas}: \citeonline{eloundou2024}
avaliam, atividade a atividade, se o acesso a um modelo de linguagem reduziria
em pelo menos metade o tempo necessário para executá-la, mantendo a qualidade do
resultado, e agregam esse veredito em uma fração de tarefas expostas. Mede,
portanto, quanto do que a ocupação faz poderia ser feito mais rapidamente com o
auxílio da ferramenta. As duas recorrem à mesma fonte ocupacional americana
(o O*NET), mas a dimensões distintas dela: o AIOE parte do inventário de
habilidades associadas a cada ocupação, ao passo que a medida de
\citeonline{eloundou2024} parte das atividades de trabalho detalhadas e das
tarefas. Respondem, por isso, a perguntas diferentes, e nada garante que ordenem
as ocupações da mesma maneira.

A escolha entre elas condiciona a interpretação dos coeficientes. O AIOE é
anterior à difusão dos modelos generativos e capta exposição à inteligência
artificial em sentido amplo, o que o torna adequado a um desenho cujo marco
temporal é o lançamento público de uma ferramenta específica apenas sob a
hipótese de que a ordenação das ocupações por exposição ampla se aproxima da
ordenação por exposição generativa; a medida de \citeonline{eloundou2024} é
contemporânea ao choque e dispensa essa hipótese, mas é específica a modelos de
linguagem e, por construção, desconsidera aplicações de inteligência artificial
que não passem por texto. Este trabalho adota o AIOE como tratamento e registra
a substituição pela medida alternativa como primeira robustez prevista na
Seção~\ref{sec:exposicoes}. Duas ressalvas, contudo, atravessam as duas famílias
e valem para qualquer leitura do Capítulo~\ref{sec:resultados}. A primeira é que
ambas medem exposição \emph{potencial}, e não adoção: o coeficiente estimado é um
gradiente sobre um índice ocupacional, não o efeito de usar a ferramenta. A
segunda é que ambas nascem na classificação ocupacional americana e precisam ser
transportadas para a classificação brasileira, operação que não é neutra e cujas
limitações de resolução ficam declaradas elo a elo na
Seção~\ref{sec:exposicao}.

\subsection{Mobilidade ocupacional, formalidade e mercado de trabalho brasileiro}

A literatura de exposição à inteligência artificial foi construída em mercados
de trabalho cuja margem de ajuste principal é o vínculo formal, e o traço que
mais distingue o caso brasileiro é a informalidade. \citeonline{ulyssea2020}
sistematiza a evidência que trata a informalidade não como resíduo nem como sala
de espera do setor formal, mas como fenômeno heterogêneo, resultante da
interação entre decisões de firmas e de trabalhadores, diferenças de
produtividade, custos de formalização e intensidade da fiscalização. Não há,
portanto, um mecanismo único a explicá-la, e sim margens distintas, que
respondem de forma distinta a cada política. \citeonline{meghir2015}, estimando
um modelo de busca e emparelhamento com microdados brasileiros, mostram que os
dois setores competem pelos mesmos trabalhadores, de modo que a fronteira entre
eles é
atravessada com frequência ao longo da vida laboral e não separa dois mercados
estanques. Daí a decisão, adotada no Capítulo~\ref{sec:metodo}, de tratar a
transição de formal para informal como desfecho próprio, com a formalidade
classificada a partir do dicionário oficial da pesquisa (carteira assinada,
vínculo estatutário ou militar e inscrição em CNPJ, no caso de conta-própria e
empregadores), em vez de tomá-la como característica fixa do trabalhador.

A segunda referência de comparação é a magnitude da rotatividade ocupacional já
existente. \citeonline{higano2023}, com dados da Pesquisa Mensal de Emprego para
as regiões metropolitanas entre 2002 e 2016, encontram que a mobilidade
ocupacional não afeta os rendimentos no agregado e que trabalhar no setor
informal aumenta a probabilidade de trocar de ocupação, dois resultados que
importam aqui, porque impedem tanto a leitura de que mudar de ocupação seja, em
si, perda, quanto a de que a troca observada seja fenômeno raro. Convém
registrar que a Pesquisa Mensal de Emprego foi encerrada em 2016 e substituída,
com metodologia atualizada e abrangência nacional, justamente pela PNAD Contínua
empregada neste trabalho: a comparação com esses resultados é indicativa, e não
direta, porque mudam a base, a cobertura geográfica e o desenho amostral.
\citeonline{wroblevski2024}, com microdados da PNAD Contínua entre 2016 e 2023,
estimam matrizes de transição e confirmam mobilidade alta dentro da
informalidade e passagem limitada para o emprego formal, com barreiras
concentradas entre mulheres, pessoas não brancas e trabalhadores de menor
escolaridade. Explorar o painel da pesquisa para esse tipo de exercício exige,
contudo, identificar longitudinalmente os indivíduos entre entrevistas
consecutivas, problema tratado por \citeonline{osorio2022} e enfrentado neste
trabalho na Seção~\ref{sec:par}. O conjunto dessa evidência fixa o contrafactual de leitura
do Capítulo~\ref{sec:resultados}: em um trimestre brasileiro comum já há muita
troca de ocupação e muita travessia da fronteira de formalidade, de modo que um
gradiente atribuído à exposição à inteligência artificial há de ser lido como um desvio sobre um pano de fundo de rotatividade elevada, e nunca como a
rotatividade em si.
